\section{Return, memory and universality of the dodecaphase cascade}
\label{sec:hmt-feig-interpretacion}

\subsection{The generating structure and the critical reading}

The Feigenbaum construction in Triadic Modular Holography links an
oriented arithmetic organization to an asymptotic scaling law.
Its starting point is the effective composition of APP, TRIT and TPK. APP evaluates
sum and product on the two sheets of the digit lattice and preserves the
quotients and remainders of those operations. TRIT determines the local regime and
the orientation of its reading. TPK selects and transports the states,
updates their records and prolongs their routes with carry, boundary and memory.
The joint discrete structure of the continuum provides the common residence
of those states and of the readers acting on them.

The prolongation
\[
w_6\longrightarrow w_{12}\longrightarrow w_{18}
\longrightarrow w_{24}\longrightarrow w_{30}
\longrightarrow R_{36}\longrightarrow G_9
\]
organizes this persistence. The initial seed is lifted while preserving its
prefixes; successive lifts change the regime without losing orientation
or memory; the coinductive boundary allows the construction to continue, and the
nine charts bring together its local readings. Nonadic holonomy expresses
phase return accompanied by an increase in memory. Thus, a visible
turn preserves its position within the complete history. From the outset,
the return operation acquires two inseparable components:
reiteration of a rule and preservation of its provenance.

The critical reader fixed on the oriented dodecaphase wheel turns this
structure into twelve phases with a uniform trace. Their representatives are
\[
\phi_m=\frac{(3m-1)\pi_{\mathrm{HMT}}}{18},
\qquad w_m=\frac1{12},\qquad 1\le m\le12.
\]
The choice of oriented representative matters because the second moment
uses its magnitude. The angular quotient records the phase, while the
representative retains the information needed to normalize it. Coordinate
\(\pi_{\mathrm{HMT}}\) comes from the coinductive readers of the same
APP--TRIT--TPK core. Within its correlated image, coordinates
\(\pi,\varphi,e\) and the dodecaphase closure of \(\alpha_{\mathrm{HMT}}\)
retain their constructive order and their own domains. The critical reader
considered here uses the angular coordinate already generated.

Normalization is determined through an integer sum:
\[
\sum_{m=1}^{12}(3m-1)^2=5394=6\cdot899,
\qquad k_m=\frac{2(3m-1)}{\sqrt{899}}.
\]
The common factor \(\pi_{\mathrm{HMT}}\) cancels when the second moment is fixed.
This cancellation preserves the provenance of the phases and produces a
particularly precise representation of their normalized reading. We obtain
\[
u_0(x)=\frac1{12}\sum_{m=1}^{12}\bigl(1-\cos(k_mx)\bigr),
\qquad f_\lambda(x)=1-\lambda u_0(x),
\qquad \sum_mw_mk_m^2=2.
\]
The profile is even and entire; it satisfies \(u_0(0)=u_0'(0)=0\) and
\(u_0''(0)=2\). For \(0<\lambda\le2/u_0(1)\), the family maps
\([-1,1]\) into itself and has a quadratic critical maximum with
\(f_\lambda''(0)=-2\lambda\). Parameter \(\lambda\) ranges over this family
after the reader, orientation, weights and scale have been fixed. The universal
number recognized at the end is thus separated from the selection
of the coefficients.

\subsection{Twelve phases, moments and spectral representation}

The profile admits a second expression bringing its phases together in a finite
operator. The nodes
\[
s_m=k_m^2=\frac{4(3m-1)^2}{899}
\]
determine the positive measure \(\mu_0=12^{-1}\sum_m\delta_{s_m}\). Its
moments \(M_r=\int s^r\,d\mu_0(s)\) define strictly positive Hankel forms
up to rank twelve. This positivity follows from evaluating
a polynomial at twelve distinct nodes: a polynomial of degree less than
twelve retains at least one nonzero evaluation. At degree twelve, the product
of the twelve nodal factors annihilates the norm and determines the reader’s
rank exactly.

Orthogonalization of these moments produces a Jacobi matrix
\(J_{12}\), real symmetric, with positive subdiagonal entries. Its functional calculus
recovers the complete profile:
\[
u_0(x)=e_0^{\mathsf T}
\bigl[I-\cos(x\sqrt{J_{12}})\bigr]e_0.
\]
The cosine sum and the Jacobi operator are two realizations of the same
reader. The moments retain all frequencies and their weights; the
orientation, route and memory marks remain in the state that
generates them. Rank twelve identifies the spectral dimension of this
specific reading. The depth of its prolongations belongs to another dimension
of the construction: the length of compatible histories and the resolution
at which they are evaluated.

The analytic geometry is also made explicit. For the weighted
deformation \(u_\eta\), in the strip \(|\eta|\le1/40\), the weights remain
positive and the second moment remains normalized. In
\(|\operatorname{Re}z|<\pi/3\) there exists an odd, univalent holomorphic function
\(b_\eta\), normalized by \(b_\eta'(0)=1\), such that
\(u_\eta=b_\eta^2\). The wheel therefore produces a conformal deformation
of the quadratic fold. Its dynamical conjugacy preserves that deformation
through \(b_\eta(1-\lambda w^2)\). The parameter-scaling theorem
developed here concerns the uniform member \(\eta=0\); the
structural properties of the weighted strip and the local persistence of
its roots retain their specific quantifiers.

\subsection{Reading depth and recovery of histories}

Nonadic depth refines the reading of a parameter, whereas dynamical
doubling changes the return horizon. Refinement with
carry expresses the first operation through
\[
L_c(x,y)=(Bx-c,By+c),\qquad
C_m=\sum_{j=1}^m c_jB^{m-j},
\]
\[
(x_m,y_m)=(B^mx_0-C_m,B^my_0+C_m).
\]
The outer support grows as \(B^m\), and the normalized record cylinder
has diameter \(B^{-m}\). The integer pair, depth and boundary
mark preserve reconstruction of quotients, remainders and prefixes.
Outer growth and inner resolution thus describe compatible
operations on the same history.

The exact composition
\(L_D^{[b]}L_C^{[a]}=L_{B^bC+D}^{[a+b]}\) allows a history to be grouped into
blocks without altering its data. A doubling return combines two
transitions; the nonadic calendar combines nine. Both groupings are
compared over a common horizon of \(9\,2^N\) events, preserving
the order of the subblocks and their carries.

The contribution of memory can be observed directly in the profile’s
inverse branches. If \(v=(u_0|_{[0,1]})^{-1}\), the two branches are
\[
\beta_\sigma(y)=\sigma v\!\left(\frac{1-y}{\lambda}\right),
\qquad \sigma\in\{-1,+1\}.
\]
The final value, together with the last sheet, reconstructs the preceding value. A
sequence of sheets allows all intermediate steps to be recovered. The branch
word accompanies the APP--TRIT--TPK records of route, boundary, carry and
memory; each record retains its function. This concrete recovery
explains what information is required to invert the critical reading.

The resolution of each return also has an explicit bound. With
\(t=\lambda u_\eta(1)/2\), \(z=(x+1)/2\) and
\[
F_{t,\eta}(z)=1-t\frac{u_\eta(2z-1)}{u_\eta(1)},
\]
one obtains
\[
\left|F_{t,\eta}^{q}(1/2)-F_{s,\eta}^{q}(1/2)\right|
\le S_q|t-s|,\qquad S_q=\sum_{j=0}^{q-1}(6\sqrt2)^j<9^q.
\]
A base-nine cylinder of depth \(m=2^N+r\) therefore allows the
return of horizon \(2^N\) to be read with diameter less than \(9^{-r}\). This
estimate gives quantitative content to arbitrary depth: for every
requested horizon and precision, it determines a sufficient depth.
Evaluation of the analytic functions and rounding have their own
remainders, which propagate together with the parameter uncertainty.

Vacancies enter the capacity balance of the record.
Comparison between blocks of 729 and 1000 possibilities distinguishes accumulated
capacity from chronological length, preserving the steps at which the former
remains constant. This accounting determines the space required to
preserve a reading and its depth. The admissibility of histories
follows, in turn, from the survival rules. Both operations
contribute to the construction and retain distinct mathematical objects.

\subsection{Nonadic chronology and doubling combinatorics}

The critical doubling word has an expression at every depth:
\[
\tau_j=(-1)^{v_2(j)},\qquad
\tau_{2j}=-\tau_j,\qquad \tau_{2j+1}=+1.
\]
Its prefixes satisfy \(W_{n+1}=W_n(-1)^nW_n\). The rule determines how
the two copies of the return are oriented and where their separation appears.
The nonadic reading of chronology preserves the complete integer
\(K=\rho_9(K)+9q_9(K)\). Reducing it modulo \(2^n\), the holonomy
that advances nine events induces translation by nine on that cycle. Since
nine is odd, this translation visits all its positions. The
permutation \(j\mapsto9j\pmod{2^n}\), which conjugates it to unit advance,
preserves the dyadic valuation of each nonzero position. The chronological remainder and quotient
jointly support that compatibility.

Bilateral memory also has an operator realization compatible
with refinement. On cycles of length \(2^n\), the shift operator
\(C_n\) is incorporated into
\[
\mathbb U_n=P_0\otimes I+P_1\otimes C_n,
\]
with the complementary projectors of the nonadic reader. In the limit, its
visible and complementary components satisfy
\[
\|T_{\infty,a}v\|^2+\|D_{\infty,a}v\|^2=\|v\|^2,
\qquad
v=T_{\infty,a}^*T_{\infty,a}v+D_{\infty,a}^*D_{\infty,a}v.
\]
These identities describe preservation and recovery within the declared space
and representation. Contraction of the renormalized maps,
which will appear later, measures another operation: the approach of their profiles
to a stable form. The history preserved by the lift and the
convergence of a functional reading are articulated at different levels.

\subsection{The ordered chart and global root selection}

The Archimedean readout of the HMT continuum transports the operations
through an isomorphism of complete ordered fields, within the declared
universe \(\mathfrak G\):
\[
\iota:\mathbb R_{\mathrm{HMT}}^{\mathfrak G}
\longrightarrow\mathbb R^{\mathfrak G}.
\]
The internal construction of the cosine through its series, the positive square root of 899,
the profile and its iterates precedes zero selection. Preservation
of sum, product, order and convergent limits gives
\[
\iota\bigl(S_n^{\mathrm H}(a)\bigr)=S_n\bigl(\iota(a)\bigr),
\qquad S_n(\lambda)=f_\lambda^{2^n}(0).
\]
Surjectivity covers all roots in the domain of that chart; order
preserves preceding returns, exact period and the condition of
lying to the right of the preceding root. In this way,
the minimum of each candidate set is also transported. The genealogy reaches
the effective parameter selection, with the second projection and its
marks preserved on the same antecedent.

The selector is fixed by \(\lambda_1=1/u_0(1)\) and by the first
new root of exact critical period \(2^n\) lying to the right of
\(\lambda_{n-1}\). The itinerary classification proves that each selected
root has word \(W_n0\). Restrictive returns halve its
period while preserving unimodal order, and induction restores
the complete word. Analyticity of returns that are not identically zero
isolates their zeros in each compact set;
comparison of signs before the first infinite-itinerary parameter
guarantees a new root at every level.

Let \(\Lambda_\tau\) be the nonempty compact set of parameters realizing
the infinite word, and let \(a_*\) be its minimum. It is proved that
\[
\lambda_n\nearrow a_*.
\]
Passage to the limit preserves each sign through a uniform bound at a fixed
horizon. If \(c_p=f_\lambda^p(0)\), the opposite signs of \(c_p\) and
\(c_{2p}\), together with \((f_\lambda^p)'(0)=0\), imply
\[
|c_p|>\frac2{B_p},\qquad
B_p=16^p\frac{16^p-1}{15}.
\]
Thus, the prefixes of increasing periods preserve the infinite word as
they accumulate. The proof uses the original sequence of first roots from
its global definition.

\subsection{The stable sheet and identification of the selected limit}

The normalized return
\[
(Rf)(x)=-\frac{f(f(-ax))}{a},\qquad a=-f(1)>0,
\]
combines two transitions, changes the scale and preserves the orientation needed
to restore a maximum of value one. The restrictive domains establish
its interpretation as a real return. The derivative of this operator includes
the variation of scale \(a\); that contribution determines the parameter
sensitivity of the complete transformation.

In the chart \(f(x)=1-x^2V((x^2-1)/(5/2))\), with coordinates \(u=10v_0\)
and \(\nu=(v_1,v_2,\ldots)\), the fourth return of the family enters a
quantified neighborhood of the fixed point \(g\). The curve crosses
its stable sheet transversely at a unique local parameter
\[
1.874038<\lambda_*<1.874039,
\]
and satisfies
\[
\|R^{4+n}f_{\lambda_*}-g\|_{\mathcal A}
<0.001666(0.83996)^n.
\]
The stable sheet is a graph of slope at most \(0.16\). The
parameter secants belong to a transverse expanding cone. Both
properties distinguish the decrease in stable deformations from the
growth of separation between parameters.

Identifying \(a_*\) with \(\lambda_*\) brings this local geometry together
with the global minimum. First, the interval between the eighth root
and the lower endpoint of the local box is excluded. Then the order
\(a_*\le\lambda_*\) allows a hypothetical separation to be studied through
the negative cone. While its transverse coordinate has absolute value less
than \(0.006\), the secant remains in the analytic domain and grows by a
factor between \(4.4034\) and \(8.0887\). A persistent separation
would therefore produce a first exit with amplitude between
\(0.006\) and \(0.0485322\).

The stable-orbit residual retains information more precise than its
norm: its components satisfy \(|e_u|\le0.16\|e_\nu\|_1\). This
relation allows their effects on the critical orbit to be bounded jointly.
The cover of the entire first-exit region shows in each case a
sign incompatible with \(\tau\), or a contracting return forcing two
successive dyadic critical values to have the same sign. The infinite
word requires opposite signs. The contradiction identifies the parameters:
\[
\lambda_n\nearrow a_*=\lambda_*.
\]
The proof combines an induction valid at every depth with a
uniform finite exclusion of the region that any separation
would reach. Its infinite scope rests on that composition.

\subsection{Universal scaling and the meaning of precision}

The degree-two holomorphic realization of the returns and the established
transversality allow their hybrid classes to be compared. For each sufficiently
high period, the canonical class has a unique local intersection
with the family. The global first roots already converge to the crossing and
preserve the doubling combinatorics; hence they eventually coincide
with the local centers when their physical periods are matched. The proof
transports the original selection to the universal metric law.

There exist \(C>0\), \(M<\infty\) and \(0<\theta<1\) such that
\[
\lambda_*-\lambda_n=C\delta_{\mathrm F}^{-n}(1+e_n),
\qquad |e_n|\le M\theta^n,
\]
and, consequently,
\[
\lim_{n\to\infty}
\frac{\lambda_{n-1}-\lambda_{n-2}}
{\lambda_n-\lambda_{n-1}}=\delta_{\mathrm F}.
\]
The factor \(\delta_{\mathrm F}\) is the expanding eigenvalue of the doubling operator
recognized at the end of the construction. The specific family preserves its
twelve phases, moments and provenance history; the limit of its
parameter ratios belongs to the same universality class of quadratic
criticality. This relation jointly explains the specificity of the
generator and the universality of the output.

Precision has three residences here. The root intervals
control their finite evaluation. The estimate
\(0.001666(0.83996)^n\) controls operator convergence on the stable
sheet. The term \(M\theta^n\) controls the transfer of that geometry to
the parameter scale through transverse holonomy. Each estimate
corresponds to its variable and operation. The eighth quotient is enclosed
around \(4.669212189150204154556096215889663928\), with width less
than \(5.808\cdot10^{-37}\); that width measures the evaluation of \(\delta_{\mathrm F,8}\).
The distance from \(\delta_{\mathrm F,8}\) to the limit requires quantifying the constants
of the parameter remainder, whose existence is already proved.

The notation keeps separate \(\delta_{\mathrm F}\), the Feigenbaum parameter-scaling
factor; \(\alpha_{\mathrm F}\), the spatial constant associated with its
renormalization problem; and \(\alpha_{\mathrm{HMT}}\), the corpus’s fine-structure
constant. The present result identifies \(\delta_{\mathrm F}\)
through the first roots of the uniform sector \(\eta=0\). The weighted
strip retains its structural and existence theorems, and its quantified
metric extension constitutes a distinct specific statement.

The contribution of the chain lies in continuity between generation,
reading, memory, selection and limit. Oriented arithmetic determines the
profile; the spectral representation preserves its twelve components; the
cylinders allow it to be resolved with arbitrary precision at each horizon;
history transport restores its paths; the ordered chart
preserves all roots and their minima; classification prolongs the
selector; and transverse expansion, together with stable contraction,
determines its asymptotic law. Universality emerges as a proved
property of this complete construction and its analytic realization.
